\begin{helper}
Let \(D\) be a digraph on a finite vertex set \(W\), and let
\(3\le\ell\).  Put
\[
  \Omega=\{\phi:\ [\ell-1]\to([n]\to W)\}.
\]
Call \(\phi\in\Omega\) bad if the coloring
\(\noderef{RandomHomomorphismColoring}(D,\ell,n,\phi)\) has a final-color
monochromatic clique of size \(s\), that is, if there is a set
\(S\subseteq[n]\) with \(|S|=s\) such that every edge inside \(S\) has color
\(\ell-1\).  If the number of bad \(\phi\in\Omega\) is strictly smaller than
\(|\Omega|\), then there exists \(\phi\in\Omega\) which is not bad.
\end{helper}
\begin{proof}
Assume, toward contradiction, that every \(\phi\in\Omega\) is bad.  Then the
map
\[
  \phi\longmapsto (\phi,\text{a proof that }\phi\text{ is bad})
\]
embeds \(\Omega\) into the subtype of bad elements of \(\Omega\).  Therefore
\[
  |\Omega|\le |\{\phi\in\Omega:\phi\text{ is bad}\}|.
\]
This contradicts the assumed strict inequality
\[
  |\{\phi\in\Omega:\phi\text{ is bad}\}|<|\Omega|.
\]
Hence at least one \(\phi\in\Omega\) is not bad.
\end{proof}
